% ============================================================
% §1  Ontologia TGP — materia jako źródło przestrzeni
% ============================================================
\section{Ontologia}\label{sec:ontologia}

\subsection{Centralne za\l{}o\.zenie}

Teoria Generowanej Przestrzeni opiera si\k{e} na jednym fundamentalnym
twierdzeniu ontologicznym:

\begin{axiom}[Przestrze\'n z~materii]\label{ax:przestrzen}
Przestrze\'n nie jest bytem niezale\.znym od materii.
\textbf{Obecno\'s\'c przestrzeni wynika z~obecno\'sci materii.}
Ka\.zda forma masy i~energii generuje lokalny wk\l{}ad do wsp\'olnego
pola $\Phi$, kt\'ore mierzy ,,g\k{e}sto\'s\'c wygenerowanej przestrzeni''
w~danym miejscu. Bez materii nie ma przestrzeni; bez przestrzeni
nie ma materii. Obie wielko\'sci s\k{a} wsp\'o\l{}zale\.zne.
\end{axiom}

\begin{remark}[Warstwa pre-substratowa]\label{rem:pre-substrat}
Aksjomat~\ref{ax:przestrzen} jest fundamentalny \emph{fizycznie},
lecz nie \emph{logicznie}: istnieje warstwa wcze\'sniejsza ni\.z
substrat $\Gamma$, w~kt\'orej promocja symetrii $\Zp$ do rangi
\emph{aktu ustanowienia r\'o\.znicy} zamyka \l{}a\'ncuch przyczynowy
od nico\'sci do rzeczywisto\'sci. Formalnie jest to
\textbf{Aksjomat~0} (dod.~\ref{app:aksjomatyka-roznicy}), z~kt\'orego
wynikaj\k{a} jako twierdzenia: niestabilno\'s\'c $\Nzero$
(tw.~\ref{thm:P3}), $\Zp$-niezmienniczo\'s\'c $H_\Gamma$
(aks.~\ref{ax:Z2-konstytutywny}) oraz nieujemno\'s\'c $\Phi \geq 0$.
Niniejsza sekcja zachowuje tradycyjn\k{a} prezentacj\k{e}
,,od substratu'' dla przejrzysto\'sci fizycznej.
\end{remark}

To odwraca standardowy obraz fizyki, w~kt\'orym przestrze\'n (lub
czasoprzestrzenie) istnieje jako uprzednio dana arena, a~materia
jedynie ,,porusza si\k{e}'' po niej. Odr\'o\.znia to r\'ownie\.z TGP od
program\'ow, w~kt\'orych przestrze\'n emerguje z~informacji lub spl\k{a}tania
(AdS/CFT, ER\,=\,EPR) --- w~TGP przestrze\'n emerguje z~\emph{materii}.

\begin{remark}[Czym nie jest TGP]
TGP \textbf{nie jest} teori\k{a} skalarno-tensorow\k{a} z~dodatkowym
polem na gotowej rozmaito\'sci. Pole $\Phi$ nie ,,mieszka'' w~uprzednio
istniej\k{a}cej czasoprzestrzeni --- ono t\k{e} czasoprzestrze\'n
\emph{konstytuuje}. Formalizm skalarno-tensorowy mo\.ze by\'c jedynie
przybli\.zeniem p\'o\'znej, s\l{}abopolowej granicy TGP, nie jej fundamentem.
\end{remark}

\begin{remark}[Tr\'ojpoziomowa ontologia TGP]\label{rem:ontologia}
TGP przyjmuje nast\k{e}puj\k{a}c\k{a} hierarchi\k{e}:
\begin{description}[style=nextline,nosep]
  \item[Poziom 0: Substrat relacyjny $\Gamma = (V, E)$]
    Fundamentalny, nie-emergentny. Posiada relacje
    s\k{a}siedztwa i~konektywno\'s\'c, ale nie posiada
    metryki. Istnieje zar\'owno w~$\mathcal{S}_0$,
    jak i~w~$\mathcal{S}_1$.
  \item[Poziom 1: Pole przestrzenno\'sci $\Phi$
    (parametr fazy metrycznej)]
    $\Phi$ \textbf{nie jest} przestrzeni\k{a} ani geometri\k{a}
    --- jest \emph{parametrem fazy metrycznej substratu},
    mierz\k{a}cym stopie\'n, w~jakim substrat wspiera relacje
    metryczne (odleg\l{}o\'sci, propagacj\k{e}).
    $\Phi = 0$: substrat istnieje, ale
    nie ma metryki (faza niemetryczna).
    $\Phi > 0$: faza metryczna --- materia generuje
    zdolno\'s\'c substratu do wspierania geometrii.
  \item[Poziom 2: Geometria efektywna $g_{\mu\nu}$]
    Emergentna w~tym sensie, \.ze wynika
    z~grubozarnienia poziomu~1. Metryka, krzywizna,
    przyczynowo\'s\'c s\k{a} w\l{}asno\'sciami
    konfiguracji~$\Phi$, nie oddzielnymi bytami.
\end{description}
Substrat jest fundamentalny i~nie-emergentny. Geometria jest
stanem tego substratu, kontrolowanym przez $\Phi$. Przej\'scie
$\Phi = 0 \to \Phi > 0$ nie tworzy przestrzeni z~nico\'sci ---
,,w\l{}\k{a}cza'' metryczn\k{a} odpowied\'z ju\.z
istniej\k{a}cego substratu. Analogia: zamarzanie wody nie
tworzy moleku\l{} --- zmienia ich faz\k{e}.
\end{remark}

\begin{remark}[Jeden substrat --- wiele konsekwencji]\label{rem:jeden-substrat}
Tr\'ojpoziomowa hierarchia powy\.zej wynika z~\textbf{jednego}
fundamentalnego obiektu: substratu relacyjnego $\Gamma = (V,E)$
z~hamiltonianem $H_\Gamma$ o~chiralnej symetrii $\mathbb{Z}_2$
(ssec.~\ref{ssec:chiralna}). Z~tego jednego obiektu
przez r\'o\.zne projekcje coarse-grainingu wynikaj\k{a}:

\begin{enumerate}[nosep]
  \item \textbf{Pole $\Phi(x)$} --- parametr fazy metrycznej,
    z~grubowania amplitudy: $\Phi \propto \langle\hat{s}^2\rangle_{\mathcal{B}(x)}$.
    Stanowi fundament \emph{ca\l{}ego} sektora grawitacyjnego
    i~kosmologicznego (sekcje~\ref{sec:pole}--\ref{sec:formalizm}).
  \item \textbf{Pole $\sigma_{ab}(x)$} --- tensor korelacji przestrzennych:
    $\sigma_{ab} \propto \langle\hat{s}_i\,\hat{s}_{i+\hat{e}_a}\rangle_{\mathcal{B}}^{\rm TF}$.
    Niesie spin-2, opisuje tensorowe mody fal grawitacyjnych
    (\S\ref{ssec:tensor-substrate}).
    Wynika z~\textbf{tego samego} $H_\Gamma$ co $\Phi$ ---
    nie jest dodatkowym postulatem.
\end{enumerate}

\noindent\textbf{Czym to \emph{nie} jest.}
Pola $\Phi$ i~$\sigma_{ab}$ nie s\k{a} polami ,,na'' przestrzeni ---
one \emph{konstytuuj\k{a}} przestrze\'n (poziom~0~$\to$~1~$\to$~2).
Nie jest to teoria skalarno-tensorowa: w~teorii skalarno-tensorowej
dwa pola s\k{a} \emph{postulowane} na gotowej rozmaito\'sci.
Tu oba emerguj\k{a} z~jednego substratu jako r\'o\.zne projekcje.

\medskip
\noindent\textbf{Sektor cechowania (program odrębny).}
Sekcja~\ref{sec:cechowanie} pokazuje, \.ze \emph{je\.zeli} substrat
posiada dodatkowy stopie\'n swobody fazowego
($\hat{s}_i \to \psi_i = |\psi_i|\,e^{i\theta_i}$),
to z~gradienta fazy emerguje pole cechowania $A_\mu$.
Jest to \textbf{oddzielny program badawczy} (nie zamkni\k{e}ty modu\l{}),
logicznie niezale\.zny od sektora grawitacyjnego.
Grawitacja w~TGP nie potrzebuje tego rozszerzenia.
\end{remark}

% ----------------------------------------------------------
\subsection{Nico\'s\'c $\Nzero$ i~funkcja stanu $\Ffield$}\label{ssec:N0}

Punktem wyj\'scia jest formalne zdefiniowanie stanu,
w~kt\'orym nie ma ani materii, ani przestrzeni.

\begin{axiom}[Nico\'s\'c]\label{ax:N0}
Fundamentalnym stanem odniesienia jest $\Nzero$ ---
absolutna nico\'s\'c. Formalnie:
\begin{equation}
    \Nzero := \bigl\{\emptyset,\;\Ffield\bigr\},
    \qquad
    \Ffield\big|_{\Nzero} = 0,
    \qquad
    \Phi\big|_{\Nzero} = 0.
\end{equation}
W~stanie $\Nzero$ nie ma przestrzeni ($\Phi = 0$), nie ma materii,
nie ma czasu, nie ma energii. Jedyn\k{a} w\l{}asno\'sci\k{a} logiczn\k{a}
jest zdefiniowanie funkcji stanu $\Ffield$ r\'ownej zeru.
\end{axiom}

$\Nzero$ nie jest pust\k{a} przestrzeni\k{a} --- jest brakiem
przestrzeni. Nie jest pr\'o\.zni\k{a} kwantow\k{a} --- jest brakiem
jakiejkolwiek struktury. Jedyno\'s\'c tego stanu (w~obr\k{e}bie
przyj\k{e}tego modelu substratu) jest wykazana
w~twierdzeniu~\ref{thm:N0-unique}.

\medskip\noindent\textbf{Precyzacja ontologiczna $\Ffield$
wzgl\k{e}dem $\Nzero$.}\;
$\Ffield$ nie jest ,,sk\l{}adnikiem ontologicznym'' nico\'sci ---
nie ,,mieszka'' w~$\Nzero$ jako cz\k{e}\'s\'c jego struktury.
$\Ffield$ jest \textbf{zewn\k{e}trznym operatorem klasyfikuj\k{a}cym}
stan: narz\k{e}dziem opisu, kt\'ore przyporz\k{a}dkowuje stanowi
$\sigma$ warto\'s\'c $\Ffield[\sigma]$.
Zapis $\Nzero = \{\emptyset, \Ffield\}$ nale\.zy czyta\'c jako:
,,stan $\Nzero$ jest charakteryzowany przez warunek $\Ffield = 0$'',
nie jako: ,,$\Nzero$ \emph{zawiera} $\Ffield$''.
Analogia: temperatura jest operatorem opisuj\k{a}cym stan gazu,
ale nie jest sk\l{}adnikiem gazu.

\begin{remark}[Substrat w~$\Nzero$: precyzacja terminologiczna]%
\label{rem:N0-substrat-precyzacja}
W~TGP wyra\.zenie \emph{absolutna nico\'s\'c} oznacza
\textbf{brak wszelkiej fizycznej struktury metrycznej}:
braku przestrzeni, materii, czasu i~propagacji.
\emph{Nie} oznacza braku substratu~$\Gamma$.

Substrat relacyjny $\Gamma = (V,E)$ jest bytem \textbf{poziomu~0}
(uwaga~\ref{rem:ontologia} i~\ref{rem:jeden-substrat}): istnieje
zar\'owno w~sektorze $\mathcal{S}_0$, jak i~w~$\mathcal{S}_1$.
Stan $\Nzero$ to faza niemetryczna substratu ---
stan, w~kt\'orym symetria $\mathbb{Z}_2$ jest nienaruszona
($\langle\hat{s}_i\rangle = 0$ dla ka\.zdego $i$),
co implikuje $\Phi = 0$ i~brak metryki.

\begin{center}
\small
\begin{tabular}{@{}p{4.5cm}p{2.6cm}p{4.5cm}@{}}
\toprule
\textbf{Stwierdzenie} & \textbf{W~TGP} & \textbf{Uzasadnienie} \\
\midrule
,,Brak metryki w~$\Nzero$'' & Prawda & $\Phi = 0$ \\
,,Brak materii w~$\Nzero$'' & Prawda & $\langle\hat{s}_i\rangle = 0$ \\
,,Brak czasu w~$\Nzero$'' & Prawda & brak propagacji, $c \to \infty$ \\
,,Brak substratu w~$\Nzero$'' & \textbf{Nieprawda} & $\Gamma$ jest fundamentalny \\
\bottomrule
\end{tabular}
\end{center}

To rozr\'o\.znienie jest niezb\k{e}dne dla sp\'ojno\'sci:
substrat musi \emph{istnie\'c} w~$\mathcal{S}_0$, by przej\'scie
$\mathcal{S}_0 \to \mathcal{S}_1$ mia\l{}o sens jako zmiana fazy ---
nie jako kreacja z~dosłownej nico\'sci.
U\.zywamy wyra\.zenia ,,absolutna nico\'s\'c''
jako \textbf{skr\'otu j\k{e}zykowego} oznaczaj\k{a}cego
,,brak fizycznej struktury (metrycznej, materialnej, czasowej)'',
nie jako twierdzenia metafizycznego o~braku czegokolwiek.
\end{remark}

\begin{proposition}[Wyprowadzenie $\Nzero$ z~funkcji stanu $\Ffield$]\label{prop:N0-from-F}
Cho\'c $\Nzero$ jest postulowany jako aksjomat~\ref{ax:N0},
mo\.zna go \textbf{wyprowadzi\'c} z~funkcji stanu $\Ffield$ ---
odwracaj\k{a}c logik\k{e} fundamentacji. Niech:
\begin{enumerate}[nosep]
  \item[(F1)] $\Omega$ --- przestrze\'n konfiguracji substratu,
  \item[(F2)] $\Ffield\colon\Omega\to\mathbb{R}_{\ge 0}$ --- ci\k{a}g\l{}a
    funkcja stanu spe\l{}niaj\k{a}ca $\Ffield[\sigma] = 0
    \;\Leftrightarrow\; \langle\hat{s}_i\rangle_\sigma = 0\;\forall i$
    (warunek~(i) z~tw.~\ref{thm:N0-unique}),
  \item[(F3)] $\mathcal{T}\subset\Omega$ --- podzbi\'or
    \textbf{konfiguracji trywialnych}: brak wewn\k{e}trznej struktury
    metrycznej (zerowa liczba relacji, brak r\'o\.znicowania),
  \item[(F4)] $\Phi\colon\Omega\to\mathbb{R}_{\ge 0}$ --- pole
    przestrzenno\'sci, zwi\k{a}zane z~$\Ffield$ przez
    $\Ffield_{\!\mathrm{abs}} = \Phi^2/\PhiZero^2$
    (def.~\ref{def:Fabs}).
\end{enumerate}
W\'owczas:
\begin{equation}\label{eq:N0-derived}
    \Nzero \;=\; \bigl\{\sigma \in \Omega :
        \Ffield[\sigma] = 0\bigr\} \;\cap\; \mathcal{T}
    \;=\; \mathcal{S}_0 \;\cap\; \mathcal{T},
\end{equation}
a~stan $\Nzero$ posiada nast\k{e}puj\k{a}ce w\l{}asno\'sci:
\begin{enumerate}
  \item $\Nzero$ jest \textbf{niepusty}: substrat w~fazie
    niemetrycznej ($T > T_c$, $v = 0$) realizuje $\Ffield = 0$.
  \item $\Nzero$ jest \textbf{jedyny} (z~dok\l{}adno\'sci\k{a} do
    symetrii $\Zp$): wynika z~tw.~\ref{thm:N0-unique}.
  \item W~$\Nzero$: $\Phi = 0$, $c = \infty$ (formalnie),
    $\hbar = \infty$, $G = \infty$ --- brak sko\'nczonych sta\l{}ych
    fizycznych, brak metryki, brak propagacji.
  \item $\Nzero$ jest \textbf{niestabilny}: z~czterech niezale\.znych
    argument\'ow (wn.~\ref{cor:N0-quadruple}).
\end{enumerate}
Postulowanie $\Nzero$ jako aksjomatu jest zatem logicznie
\textbf{r\'ownowa\.zne} z~uznaniem, \.ze funkcja stanu $\Ffield$
posiada (co najmniej) jedno zero w~sektorze trywialnym ---
co jest gwarantowane przez faz\k{e} niemetryczn\k{a} substratu.
\end{proposition}

\begin{proof}
Dow\'od w~trzech cz\k{e}\'sciach.

\emph{Niepusto\'s\'c.}\;
W~modelu substratu (hamiltoniano~\ref{eq:H-Gamma}), dla $T > T_c$
(faza nieuporz\k{a}dkowana), $\langle\hat{s}_i\rangle = 0$ dla
ka\.zdego $i$ z~symetrii $\Zp$. Zatem $\Ffield[\sigma_{T>T_c}] = 0$
(z~warunku F2). Jednocze\'snie brak parametru porz\k{a}dku ($v = 0$)
oznacza brak struktury metrycznej ($\Phi = 0$), wi\k{e}c
$\sigma_{T>T_c} \in \mathcal{T}$. St\k{a}d
$\sigma_{T>T_c} \in \Nzero$: zbi\'or jest niepusty.

\emph{Jedyno\'s\'c.}\;
Z~twierdzenia~\ref{thm:N0-unique}: zbi\'or $\mathcal{S}_0
= \{\sigma : \Ffield[\sigma] = 0\}$ jest sp\'ojny
i~stanowi jedn\k{a} orbit\k{e} $\Zp$.
Przeci\k{e}cie z~$\mathcal{T}$ (konfiguracje trywialne)
daje jedn\k{a} klas\k{e} r\'ownowa\.zno\'sci, poniewa\.z
\textbf{ka\.zdy} element $\mathcal{S}_0$ jest trywialny:
je\'sli $\langle\hat{s}_i\rangle = 0\;\forall i$, to
$\Phi = \langle\hat{s}^2\rangle_{\mathrm{cg}} = 0$
(fluktuacje zaniedbywalnie ma\l{}e w~granicy termodynamicznej
powy\.zej $T_c$), wi\k{e}c brak metryki.
Zatem $\mathcal{S}_0 \cap \mathcal{T} = \mathcal{S}_0$:
$\Nzero$ jest jedynym stanem z~$\Ffield = 0$.

\emph{W\l{}asno\'sci fizyczne.}\;
Z~$\Phi = 0$ i~aksjomat\'ow~\ref{ax:c}--\ref{ax:G}:
$c(0) = c_0(\PhiZero/0)^{1/2} = \infty$
(formalnie: brak sko\'nczonej pr\k{e}dko\'sci propagacji,
bo nie ma metryki). Analogicznie $\hbar(0) = \infty$,
$G(0) = \infty$. Sta\l{}e fizyczne s\k{a} nieokre\'slone ---
sp\'ojne z~brakiem struktury.
Niestabilno\'s\'c: z~wn.~\ref{cor:N0-quadruple}
(termodynamika, r\'ownanie pola, degeneracja wariacyjna,
miara konfiguracyjna).
\end{proof}

\begin{remark}[Aksjomat czy twierdzenie?]\label{rem:N0-aksjomat-vs-twierdzenie}
Stwierdzenie~\ref{prop:N0-from-F} pokazuje, \.ze $\Nzero$ jest
\textbf{wyprowadzalny} z~w\l{}asno\'sci funkcji stanu $\Ffield$
i~modelu substratu. Aksjomat~\ref{ax:N0} mo\.zna zatem traktowa\'c
jako \textbf{skr\'ot poj\k{e}ciowy}: jest logicznie r\'ownowa\.zny
istnieniu fazy niemetrycznej substratu. Zachowujemy go jako aksjomat
ze wzgl\k{e}d\'ow pedagogicznych --- bezpo\'srednio wprowadza
centralny byt teorii.
\end{remark}

\begin{theorem}[Jedyno\'s\'c $\Nzero$ w~modelu substratu]\label{thm:N0-unique}
Niech $\Omega$ b\k{e}dzie przestrzeni\k{a} konfiguracji substratu
z~hamiltonianem $H_\Gamma$ o~symetrii $\Zp$, i~niech
$\Ffield\colon\Omega\to\mathbb{R}_{\ge 0}$ b\k{e}dzie funkcj\k{a}
stanu spe\l{}niaj\k{a}c\k{a}:
\begin{enumerate}[nosep]
  \item $\Ffield[\sigma] = 0 \;\Leftrightarrow\;
    \langle\hat{s}_i\rangle_\sigma = 0\;\;\forall i$
    (symetria $\Zp$ nienaruszona),
  \item $\Ffield$ jest ci\k{a}g\l{}a wzgl\k{e}dem topologii
    produktowej na $\Omega$,
  \item \'srednie $\langle\hat{s}_i\rangle_\sigma$ s\k{a}
    \textbf{afiniczne} (liniowe) w~mierze $\sigma$
    (spe\l{}nione dla \'srednich Gibbsa z~hamiltonianem
    liniowym w~$\sigma$),
  \item topologia $\Omega$ jest produktowa: $\Omega = \prod_{i\in V}[-1,1]$.
\end{enumerate}
Wtedy, \textbf{w~obr\k{e}bie powy\.zszych za\l{}o\.ze\'n},
zbi\'or $\mathcal{S}_0 = \{\sigma:\Ffield[\sigma]=0\}$
jest \textbf{sp\'ojny} i~stanowi dok\l{}adnie jedn\k{a}
klas\k{e} r\'ownowa\.zno\'sci orbitalnej symetrii $\Zp$.
W~szczeg\'olno\'sci $\Nzero$ jest jedynym stanem
funkcjonalnym (z~dok\l{}adno\'sci\k{a} do symetrii substratu),
w~kt\'orym $\Ffield = 0$.

\medskip\noindent\emph{Zastrze\.zenie.} Wynik zale\.zy
od za\l{}o\.zenia~(iii) (liniowo\'s\'c \'srednich)
i~za\l{}o\.zenia~(iv) (topologia produktowa). Os\l{}abienie
kt\'orego\'s z~nich --- np.\ przej\'scie do substratu
o~nietrywialnej topologii lub modelu ze spontanicznie
z\l{}amanymi korelacjami d\l{}ugozasi\k{e}gowymi --- mog\l{}oby
prowadzi\'c do wielosk\l{}adnikowego $\mathcal{S}_0$.
Jedyno\'s\'c $\Nzero$ jest zatem \textbf{twierdzeniem modelu},
nie absolutnym faktem metafizycznym.
\end{theorem}

\begin{proof}
Dow\'od przebiega w~trzech krokach, opieraj\k{a}c si\k{e}
na lemacie~\ref{lem:fixed-point} poni\.zej.

\emph{Krok~1 (Zamkni\k{e}to\'s\'c).}
Warunek $\langle\hat{s}_i\rangle_\sigma = 0$ dla ka\.zdego $i$
wyznacza przeci\k{e}cie hiperpowierzchni
$\mathcal{H}_i = \{\sigma\in\Omega : \langle\hat{s}_i\rangle_\sigma=0\}$.
Ka\.zda $\mathcal{H}_i$ jest zamkni\k{e}ta (przeciwobraz zera
funkcji ci\k{a}g\l{}ej), wi\k{e}c
$\mathcal{S}_0 = \bigcap_i \mathcal{H}_i$ jest zamkni\k{e}ty.

\emph{Krok~2 (Identyfikacja ze zbiorem punkt\'ow sta\l{}ych).}
Dzia\l{}anie $\Zp$: $\iota(\sigma)_i = -\sigma_i$.
Warto\'s\'c oczekiwana spe\l{}nia $\langle\hat{s}_i\rangle_{\iota(\sigma)}
= -\langle\hat{s}_i\rangle_\sigma$ (r\'ownowariantno\'s\'c).
Zatem $\langle\hat{s}_i\rangle_\sigma = 0\;\forall i$ wtedy
i~tylko wtedy, gdy $\sigma$ le\.zy w~zbiorze punkt\'ow sta\l{}ych
$\Omega^{\Zp}$ dzia\l{}ania $\Zp$ \textbf{na poziomie \'srednich}.
Z~lematu~\ref{lem:fixed-point}:
$\Omega^{\Zp}$ jest sp\'ojny (zwarty i~\l{}\k{a}kowo sp\'ojny).

\emph{Krok~3 (Jedyno\'s\'c orbity).}
Dzia\l{}anie $\Zp$ na $\mathcal{S}_0$ jest trywialne
(ka\.zdy element jest punktem sta\l{}ym). Zatem
$\mathcal{S}_0$ jest jedn\k{a} orbit\k{a} --- i~stanowi
dok\l{}adnie jedn\k{a} klas\k{e} r\'ownowa\.zno\'sci.
\end{proof}

\begin{lemma}[Sp\'ojno\'s\'c zbioru punkt\'ow sta\l{}ych]%
\label{lem:fixed-point}
Niech $\Omega = \prod_{i\in V} [-1,1]$ z~topologi\k{a} produktow\k{a},
a~$\Zp$ dzia\l{}a przez $\iota(\sigma)_i = -\sigma_i$.
Zbi\'or punkt\'ow sta\l{}ych na poziomie \'srednich makroskopowych,
\[
  \Omega^{\Zp}_{\mathrm{avg}}
  = \bigl\{\sigma\in\Omega :
    \langle\hat{s}_i\rangle_\sigma = 0\;\;\forall i\bigr\},
\]
jest niepusty, zwarty i~sp\'ojny (w~szczeg\'olno\'sci
\l{}\k{a}kowo sp\'ojny).
\end{lemma}

\begin{proof}
\emph{Niepusto\'s\'c}: konfiguracja $\sigma_i = 0\;\forall i$
(je\'sli dozwolona) lub dowolna kombinacja symetryzowana
$\mu = \frac{1}{2}(\delta_\sigma + \delta_{\iota(\sigma)})$
daje $\langle\hat{s}_i\rangle = 0$.

\emph{Zwarto\'s\'c}: $\Omega^{\Zp}_{\mathrm{avg}}$ jest
przeci\k{e}ciem zamkni\k{e}tych zbior\'ow w~zwartej
przestrzeni $\Omega$ (tw.~Tichonowa).

\emph{Sp\'ojno\'s\'c}: Niech $\sigma, \sigma' \in
\Omega^{\Zp}_{\mathrm{avg}}$. Definiujemy \'scie\.zk\k{e}
$\sigma(t) = (1-t)\sigma + t\sigma'$ dla $t\in[0,1]$.
R\'ownowariantno\'s\'c $\langle\hat{s}_i\rangle$ wzgl\k{e}dem
$\Zp$ implikuje, \.ze $\langle\hat{s}_i\rangle$ jest
\textbf{nieparzysta} w~$\sigma$. Kombinacja wypuk\l{}a dw\'och
punkt\'ow z~$\langle\hat{s}_i\rangle = 0$ daje
$\langle\hat{s}_i\rangle_{\sigma(t)} = 0$ pod warunkiem
liniowo\'sci \'sredniej w~konfiguracji (spe\l{}nione
dla \'srednich Gibbsa z~hamiltonianem liniowym w~$\sigma$)
lub --- og\'olniej --- z~tw.~Brouwera o~punkcie sta\l{}ym
zastosowanego do odwzorowania $(1-t)\mathrm{id} + t\,\iota$,
kt\'ore jest homotopijne z~identyczno\'sci\k{a} na $\Omega$.
\end{proof}

\begin{definition}[Funkcja stanu $\Ffield$]\label{def:F}
Funkcja stanu $\Ffield$ jest miar\k{a} \textbf{odej\'scia od nico\'sci}.
Przyjmuje warto\'sci:
\begin{itemize}[nosep]
  \item $\Ffield = 0$: pe\l{}na symetria, brak struktury, stan $\Nzero$.
  \item $\Ffield > 0$: lokalne z\l{}amanie symetrii --- istnieje materia
        i~wygenerowana przez ni\k{a} przestrze\'n.
\end{itemize}
$\Ffield$ mo\.zna rozumie\'c jako globaln\k{a} miar\k{e} ,,ilo\'sci
istnienia'' --- ile struktury (materii, przestrzeni, relacji)
jest obecne w~danym opisie. Dynamika $\Ffield$ jest okre\'slona
przez pr\k{a}d zachowania $j_{\Ffield}^\mu$
(stwierdzenie~\ref{prop:F-current}), a~jej ewolucja wynika
z~r\'ownania pola $\Phi$ (r\'ownanie~\ref{eq:F-evolution}).

Na poziomie formalizmu (sekcja~\ref{ssec:F-Phi}) $\Ffield$ realizuje
si\k{e} jako \textbf{absolutna funkcja stanu}
$\Ffield_{\!\mathrm{abs}}(x) = \Phi(x)^2/\PhiZero^2$,
kt\'ora jest zerowa dok\l{}adnie w~$\Nzero$ ($\Phi = 0$).
Nale\.zy j\k{a} odr\'o\.znia\'c od perturbacyjnej miary
$\Floc = (\Phi - \PhiZero)^2$, kt\'ora jest zerowa
na tle referencyjnym ($\Phi = \PhiZero$) i~s\l{}u\.zy
do rachunk\'ow wok\'o\l{} tego t\l{}a
(uwaga~\ref{rem:F-dwa-poziomy}).
\end{definition}

\begin{axiom}[Zasada zerowej sumy]\label{ax:zero}
Widziany z~zewn\k{a}trz, zamkni\k{e}ty uk\l{}ad (Wszech\'swiat) jest
nieodr\'o\.znialny od~$\Nzero$:
\begin{equation}\label{eq:zero-sum}
    \Ffield_{\mathrm{ext}}\bigl(\text{Wszech\'swiat}\bigr) = 0.
\end{equation}
Wewn\k{a}trz: $\Ffield_{\mathrm{int}}(x) > 0$ lokalnie. Ca\l{}kowita
,,ilo\'s\'c istnienia'' netto jest zerowa --- Wszech\'swiat jest
samowyzerowuj\k{a}c\k{a} si\k{e} fluktuacj\k{a} nico\'sci.
\end{axiom}

Analogia: Wszech\'swiat jest jak fala na morzu --- lokalnie jest
wzniesienie i~obni\.zenie, ale \'sredni poziom wody si\k{e} nie zmienia.

\begin{remark}[Precyzacja zerowej sumy --- dwa
  niezale\.zne warunki]\label{rem:zero-precyzacja}
Zasada zerowej sumy rozk\l{}ada si\k{e} na \textbf{dwa
logicznie niezale\.zne warunki}:

\medskip\noindent
\textbf{(ZS1) Zerowa suma chiralna ($\Zp$).}
Nale\.zy rozr\'o\.zni\'c dwa obiekty:
\begin{itemize}[nosep]
  \item $\Ffield_{\mathrm{int}}(x) \geq 0$ --- lokalna miara
        odej\'scia od symetrii \emph{wewn\k{a}trz} uk\l{}adu.
        Jest nieujemna i~niezerowa wsz\k{e}dzie, gdzie istnieje
        materia i~przestrze\'n.
  \item $\Ffield_{\mathrm{ext}}$ --- wypadkowy opis
        \emph{z~zewn\k{a}trz}: suma netto ze znakiem.
\end{itemize}
Zasada zerowej sumy (\ref{eq:zero-sum}) nie m\'owi, \.ze
$\int \Ffield_{\mathrm{int}} = 0$ (to jest niemo\.zliwe,
bo $\Ffield_{\mathrm{int}} \geq 0$). M\'owi, \.ze
\textbf{asymetria chiralna ze znakiem} $\Delta(x)$
(r\'o\.znica $\rho(s,x) - \rho(-s,x)$ w~opisie $\Zp$,
odpowiadaj\k{a}ca lokalnej warto\'sci $\langle\hat{s}\rangle$)
ca\l{}kuje si\k{e} do zera:
\begin{equation}\label{eq:zero-sum-precise}
    \int_\Sigma \Delta(x)\,\sqrt{h}\,d^3x = 0.
\end{equation}
Lokalna asymetria mo\.ze by\'c niezerowa (i~jest --- to jest
materia), ale globalna suma ze znakiem wynosi zero.
Na poziomie substratu: w~fazie z~$\Zp$ z\l{}aman\k{a}
istniej\k{a} domeny $\langle\hat{s}\rangle = +v$
i~$\langle\hat{s}\rangle = -v$; globalnie ca\l{}ka
$\int\langle\hat{s}\rangle = 0$.

\medskip\noindent
\textbf{(ZS2) Zerowa suma przestrzenna ($\Phi$).}
W~j\k{e}zyku pola $\Phi$ zasada zerowej sumy
przyjmuje posta\'c drugiego warunku
(r\'ownanie~\ref{eq:zero-constraint}):
\begin{equation}\label{eq:zero-sum-Phi}
    \int_\Sigma \bigl(\Phi(\mathbf{x}) - \PhiZero\bigr)\,
    \sqrt{h}\,d^3x = 0.
\end{equation}
\'Srednia perturbacja pola przestrzenno\'sci po ca\l{}ej
hypersurface wynosi zero: nadwy\.zka wygenerowanej
przestrzeni ($\Phi > \PhiZero$) wok\'o\l{} \'zr\'ode\l{}
jest dok\l{}adnie kompensowana deficytem
($\Phi < \PhiZero$) w~regionach mi\k{e}dzy nimi.

\begin{remark}[Autoreferencyjno\'s\'c ZS2]\label{rem:ZS2-self-ref}
Warunek~\eqref{eq:zero-sum-Phi} zawiera $\sqrt{h}$,
kt\'ore w~TGP zale\.zy od $\Phi$:
$\sqrt{h} = (\Phi/\PhiZero)^{3/2}$
(stw.~\ref{prop:conformal-unique}).
Nie jest to b\l{}\k{a}d ko\l{}owy, lecz
\textbf{autoreferencyjny warunek selekcji}:
dla danego $\Phi(\mathbf{x})$ metryka jest
okre\'slona, ca\l{}ka jest obliczalna,
a~ZS2 wybiera dopuszczalne konfiguracje.

W~przybli\.zeniu $\delta\Phi = \Phi - \PhiZero \ll \PhiZero$:
\begin{equation}\label{eq:ZS2-expansion}
  0 = \int\delta\Phi\,d^3x
    + \frac{3}{2\PhiZero}\int(\delta\Phi)^2\,d^3x
    + \mathcal{O}(\delta\Phi^3/\PhiZero^2).
\end{equation}
Do wiod\k{a}cego rz\k{e}du ZS2 jest ,,p\l{}aski'':
$\langle\delta\Phi\rangle_{\mathrm{koord}} = 0$.
Korekcja metryczna
$\langle\delta\Phi\rangle_{\mathrm{koord}}
= -\frac{3}{2\PhiZero}\langle(\delta\Phi)^2\rangle_{\mathrm{koord}}
< 0$
przesuwa \'sredni\k{a} koordynatow\k{a} w~stron\k{e} ujemn\k{a}
(regiony o~$\Phi > \PhiZero$ ,,wa\.z\k{a} wi\k{e}cej'' fizycznie).
Efekt ten jest rz\k{e}du $\delta\Phi^2/\PhiZero$
i~zaniedbywalny w~re\.zimie s\l{}abopolowym.
\end{remark}

\medskip\noindent
\textbf{Zwi\k{a}zek mi\k{e}dzy ZS1 a~ZS2 --- hierarchia.}
\begin{enumerate}[nosep]
  \item \textbf{ZS1 jest fundamentalna} --- wynika
    bezpo\'srednio z~$\Zp$-symetrii substratu i~postulatu,
    \.ze Wszech\'swiat jako ca\l{}o\'s\'c jest
    nieodr\'o\.znialny od $\Nzero$.
  \item \textbf{ZS2 jest efektywna} --- operuje na poziomie
    pola $\Phi$ po coarse-grainingu.
  \item ZS2 \textbf{nie wynika z~ZS1}: ZS1 dotyczy
    \emph{pierwszego momentu} ($\langle\hat{s}\rangle$),
    ZS2 dotyczy \emph{drugiego momentu}
    ($\Phi \propto \langle\hat{s}^2\rangle$,
    wniosek~\ref{cor:F-s}). Warunek na \'sredni\k{a}
    kwadratu nie wynika z~warunku na \'sredni\k{a} znaku.
  \item Teoria \textbf{wymaga obu}: ZS1 zapewnia globaln\k{a}
    neutralno\'s\'c chiralną, ZS2 --- bilans przestrzenny.
    S\k{a} to \textbf{dwa niezale\.zne aksjomaty}.
\end{enumerate}
\end{remark}

\begin{proposition}[Wyprowadzenie ZS2 z~dynamiki na zwartej topologii]%
\label{prop:ZS2-from-dynamics}
Je\.sli topologia przestrzenna $\Sigma$ jest zwarta (np.~$S^3$ lub $T^3$),
to ZS2 \textbf{wynika} z~r\'ownania pola TGP
(r.~\eqref{eq:full-field-alt}) i~nie stanowi niezale\.znego aksjomatu.
\end{proposition}

\begin{proof}
Ca\l{}kujemy statyczne r\'ownanie pola~\eqref{eq:full-field-alt}
po zwartej hiperpowierzchni $\Sigma$ (bez brzegu):
\begin{equation}\label{eq:ZS2-integral}
  \underbrace{\int_\Sigma \nabla^2\Phi\,d^3x}_{= 0\;\text{(Gauss, }\partial\Sigma = \varnothing\text{)}}
  \;+\; \alpha\int_\Sigma \frac{(\nabla\Phi)^2}{\Phi}\,d^3x
  \;+\; \frac{\beta}{\PhiZero}\int_\Sigma \Phi^2\,d^3x
  \;-\; \frac{\gamma}{\PhiZero^2}\int_\Sigma \Phi^3\,d^3x
  = -q\PhiZero\,M_{\mathrm{tot}},
\end{equation}
gdzie $M_{\mathrm{tot}} = \int_\Sigma \rho\,d^3x$.

Podstawiamy $\Phi = \PhiZero(1 + \varphi)$,
$|\varphi| \ll 1$ (re\.zim s\l{}abopolowy), $\beta = \gamma$:
\begin{align*}
  &\alpha\PhiZero\int_\Sigma (\nabla\varphi)^2 d^3x
  + \gamma\PhiZero \int_\Sigma
  \bigl[(1{+}\varphi)^2 - (1{+}\varphi)^3\bigr]\,d^3x
  = -q\PhiZero M_{\mathrm{tot}},\\
  &(1{+}\varphi)^2 - (1{+}\varphi)^3
  = -\varphi - 2\varphi^2 - \varphi^3.
\end{align*}
Do wiod\k{a}cego rz\k{e}du ($\varphi \ll 1$,
pomijaj\k{a}c $(\nabla\varphi)^2$ i~$\varphi^2$):
\begin{equation}\label{eq:ZS2-leading}
  -\gamma\PhiZero \int_\Sigma \varphi\,d^3x
  = -q\PhiZero\,M_{\mathrm{tot}}
  \quad\Longrightarrow\quad
  \int_\Sigma \varphi\,d^3x
  = \frac{q}{\gamma}\,M_{\mathrm{tot}}.
\end{equation}
Dla \textbf{zamkni\k{e}tego} Wszech\'swiata
ca\l{}kowita ,,masa generuj\k{a}ca''
$M_{\mathrm{tot}}^{\mathrm{gen}}$ musi znikn\k{a}\'c,
poniewa\.z w~TGP masa i~przestrze\'n
s\k{a} wz\k{a}jemnie konstytuowane
(aksjomat~\ref{ax:zrodlo}):
ka\.zde \'zr\'od\l{}o generuje $\Phi > \PhiZero$,
lecz ca\l{}kowita perturbacja musi by\'c kompensowana
(brak ,,zewn\k{e}trznego'' zasobu przestrzeni).
Formalnie: na $S^3$ twierdzenie Gaussa
dla r\'ownania Poissona wymaga zerowej ca\l{}kowitej masy.
St\k{a}d $M_{\mathrm{tot}}^{\mathrm{gen}} = 0$
i~$\int_\Sigma \varphi\,d^3x = 0$, co jest ZS2.

\medskip\noindent
\textbf{Wniosek:} dla topologii zwartej (zamkni\k{e}tej)
ZS2 jest \textbf{konsekwencj\k{a}} r\'ownania pola plus
twierdzenia Gaussa, a~nie niezale\.znym postulatem.
Dla topologii otwartej ($\mathbb{R}^3$) ZS2 pozostaje
aksjomatem wymaganym na mocy filozofii
,,Wszech\'swiat jako ca\l{}o\'s\'c = $\Nzero$''.
\end{proof}

% ----------------------------------------------------------
\subsection{Mapa poj\k{e}\'c: mikro $\to$ makro}\label{ssec:mapa}

Poni\.zsza tabela porz\k{a}dkuje relacje mi\k{e}dzy kluczowymi
wielko\'sciami na r\'o\.znych poziomach opisu TGP.

\begin{center}
\small
\begin{tabular}{@{}>{\raggedright}p{2.3cm}>{\raggedright}p{2.5cm}p{9.5cm}@{}}
\toprule
\textbf{Poziom} & \textbf{Wielko\'s\'c} & \textbf{Znaczenie i~relacje} \\
\midrule
\multicolumn{3}{@{}l}{\textbf{Mikroskopowy (substrat $\Gamma$)}} \\[2pt]
Operator & $\hat{s}_i$ & Lokalna amplituda wzbudzenia substratu \\
Asymetria chiralna & $\Delta(x) = \langle\hat{s}(x)\rangle$
  & Znakowa miara z\l{}amania $\Zp$; $\int\Delta\sqrt{h}\,d^3x = 0$ (ZS1) \\
Funkcja stanu (mikro) & $\Ffield_i = \langle\hat{s}_i\rangle^2$
  & Kwadrat asymetrii; $\Ffield = 0 \Leftrightarrow$ symetria niez\l{}amana \\
\midrule
\multicolumn{3}{@{}l}{\textbf{Po coarse-grainingu (pole $\Phi$)}} \\[2pt]
Pole przestrzenno\'sci & $\Phi(x) \propto \langle\hat{s}^2\rangle_{\mathrm{cg}}$
  & Parametr fazy metrycznej substratu;
    $\Phi = 0$: faza niemetryczna, $\Phi > 0$: faza metryczna \\
Perturbacja & $\varphi = \Phi/\PhiZero - 1$
  & Bezwymiarowe odchylenie od t\l{}a; $\int\varphi\sqrt{h}\,d^3x = 0$ (ZS2) \\
\midrule
\multicolumn{3}{@{}l}{\textbf{Globalny (funkcje stanu)}} \\[2pt]
$\Ffield_{\!\mathrm{abs}}$ & $\Phi^2/\PhiZero^2$
  & Absolutna miara odej\'scia od $\Nzero$; zerowa dok\l{}adnie w~$\Phi = 0$ \\
$\Floc$ & $(\Phi - \PhiZero)^2$
  & Perturbacyjna miara odej\'scia od t\l{}a $\PhiZero$;
    zerowa w~$\Phi = \PhiZero$ \\
\midrule
\multicolumn{3}{@{}l}{\textbf{Emergentny (geometria)}} \\[2pt]
Metryka efektywna & $g_{\mu\nu}^{\mathrm{eff}}[\Phi]$
  & Obiekt emergentny z~konfiguracji $\Phi$;
    okre\'slona tylko dla $\Phi > 0$ \\
Sta\l{}e efektywne & $c(\Phi), \hbar(\Phi), G(\Phi)$
  & Pola wynikowe, nie oddzielne postulaty \\
\bottomrule
\end{tabular}
\end{center}

\begin{remark}[Co jest pierwotne?]\label{rem:co-pierwotne}
W~hierarchii poj\k{e}\'c TGP:
\begin{enumerate}[nosep]
  \item \textbf{Pierwotne}: substrat $\Gamma$ z~operatorami
    $\hat{s}_i$ i~symetri\k{a} $\Zp$.
  \item \textbf{Pochodne (z~coarse-grainingu)}: pole $\Phi$,
    perturbacja $\varphi$, funkcje stanu $\Ffield_{\!\mathrm{abs}}$
    i~$\Floc$.
  \item \textbf{Emergentne}: metryka $g_{\mu\nu}$, sta\l{}e
    $c,\hbar,G$, geometria.
\end{enumerate}
$\Delta(x)$ i~$\varphi(x)$ \textbf{nie s\k{a} tym samym}, ale
s\k{a} \'sci\'sle powi\k{a}zane: $\Delta$ jest mikroskopow\k{a}
asymetri\k{a} chiralną, $\varphi$ jest efektywnym polem uzyskanym
po u\'srednieniu. Na poziomie pola \'sredniego: $\Phi \propto
\Delta^2$, wi\k{e}c $\varphi$ zale\.zy od $\Delta^2$, ale nie
od znaku $\Delta$.
\end{remark}

% ----------------------------------------------------------
\subsection{Status ontologiczny $\PhiZero$ i~tablica stan\'ow}\label{ssec:PhiZero-status}

\begin{remark}[Co oznacza $\PhiZero$]\label{rem:PhiZero-status}
$\PhiZero$ \textbf{nie jest stanem bez materii}.
Aksjomat~\ref{ax:przestrzen} m\'owi, \.ze bez materii nie ma
przestrzeni ($\Phi = 0$). Zatem stan $\Phi = \PhiZero > 0$
\textbf{wymaga} obecno\'sci materii --- jest to
\emph{efektywne t\l{}o ustalone przez globaln\k{a}
(\'sredni\k{a}) zawarto\'s\'c materii--energii Wszech\'swiata}.

Precyzyjnie:
\begin{itemize}[nosep]
  \item $\PhiZero$ jest \textbf{warto\'sci\k{a} referencyjn\k{a}}:
    \'sredni\k{a} warto\'sci\k{a} $\Phi$ na du\.zych skalach
    kosmologicznych, ustaloną przez ca\l{}kowit\k{a} zawarto\'s\'c
    materii--energii.
  \item $\PhiZero$ jest jednorodnym rozwi\k{a}zaniem r\'ownania
    pola (stw.~\ref{prop:vacuum-condition}) przy \'sredniej
    g\k{e}sto\'sci $\bar{\rho}$.
  \item Okre\'slenie ,,pr\'o\.znia lokalna'' ($\Phi = \PhiZero$)
    oznacza: region bez \emph{lokalnych} \'zr\'ode\l{}, ale
    z~t\l{}em wygenerowanym przez \emph{globaln\k{a}} materi\k{e}.
  \item Stan $\Phi = 0$ ($\Nzero$) jest fundamentalnie
    r\'o\.zny: to brak przestrzeni, nie ,,pusta przestrze\'n''.
\end{itemize}
Bez tego rozr\'o\.znienia mo\.zna by\l{}oby zarzuci\'c, \.ze TGP
,,najpierw twierdzi, \.ze bez materii nie ma przestrzeni,
a~potem wprowadza t\l{}o przestrzenne jak w~zwyk\l{}ej teorii
pola''. To nie jest sprzeczno\'s\'c: $\PhiZero$ \emph{jest}
generowane przez materi\k{e} --- jest po prostu \'srednim
wynikiem tej generacji, przyj\k{e}tym za punkt odniesienia.
\end{remark}

\begin{remark}[Tablica stan\'ow TGP]\label{rem:tablica-stanow}
Poni\.zsza tabela systematyzuje kluczowe stany i~poj\k{e}cia
pola $\Phi$, od absolutnej nico\'sci po silne pole:

\begin{center}
\small
\begin{tabular}{@{}>{\raggedright}p{2.5cm}cp{10cm}@{}}
\toprule
\textbf{Stan} & $\Phi$ & \textbf{Interpretacja} \\
\midrule
Nico\'s\'c $\Nzero$ & $0$
  & Brak przestrzeni, materii, czasu. Substrat istnieje,
    ale w~fazie niemetrycznej. $\Ffield_{\!\mathrm{abs}} = 0$. \\[4pt]
T\l{}o referencyjne & $\PhiZero$
  & \'Sredni stan generowany przez kosmologiczn\k{a} zawarto\'s\'c
    materii. ,,Pr\'o\.znia'' TGP: brak lokalnych \'zr\'ode\l{}, ale
    przestrze\'n istnieje (generowana globalnie).
    $\Floc = 0$, $\Ffield_{\!\mathrm{abs}} = 1$. \\[4pt]
Pr\'o\.znia lokalna & $\approx \PhiZero$
  & Region daleko od lokalnych \'zr\'ode\l{};
    $\Phi$ odchyla si\k{e} nieznacznie od $\PhiZero$.
    Odpowiada pr\'o\.zni w~sensie codziennym. \\[4pt]
Region \'zr\'od\l{}owy & $> \PhiZero$
  & Otoczenie mas: $\Phi$ podniesione przez lokalne \'zr\'od\l{}a.
    Metryka efektywna odbiega od Minkowskiego. \\[4pt]
Silne pole & $\gg \PhiZero$
  & Wn\k{e}trza gwiazd neutronowych, okolice CD.
    Sta\l{}e $c, \hbar, G$ silnie zmodyfikowane. \\[4pt]
Domena metryczna & $> 0$
  & Sp\'ojny region z~$\Phi > 0$, w~kt\'orym propagacja,
    odleg\l{}o\'sci i~geometria maj\k{a} sens. \\[4pt]
Stan zamro\.zony (CD) & $\to\infty$
  & Efektywny lokalny odpowiednik $\Nzero$:
    $c \to 0$, brak propagacji, brak dost\k{e}pu
    informacyjnego. \\
\bottomrule
\end{tabular}
\end{center}
\end{remark}

% ----------------------------------------------------------
\subsection{Chiralna symetria substratu}\label{ssec:chiralna}\label{ax:symmetry}

Na substracie relacyjnym $\Gamma = (V, E)$ ka\.zdemu
w\k{e}z\l{}owi $i \in V$ przypisujemy \textbf{operator
amplitudy} $\hat{s}_i$ --- skal\k{a} opisuj\k{a}c\k{a}
lokalny stopie\'n wzbudzenia substratu.

\begin{axiom}[Symetria $\Zp$ jako akt konstytutywny]%
\label{ax:Z2-konstytutywny}\label{prop:Z2}
Hamiltonian substratu (v2, pivot aksjomatyczny 2026-04-24; por.
\texttt{KNOWN\_ISSUES.md} w~\texttt{tgp-core-paper/}):
\begin{equation}\label{eq:H-Gamma}
    H_\Gamma = \sum_i \biggl[
      \frac{\hat{\pi}_i^2}{2\mu}
    + \frac{m_0^2}{2}\hat{s}_i^2
    + \frac{\lambda_0}{4}\hat{s}_i^4
    \biggr]
    + J\!\sum_{\langle ij\rangle} A_{ij}\,
      \hat{s}_i^2\,\hat{s}_j^2\,
      \bigl(\hat{s}_j^2 - \hat{s}_i^2\bigr)^2
\end{equation}
(z $J>0$, $\lambda_0>0$, $A_{ij}\in\{0,1\}$). Wi\k{a}zanie jest gradientowym
cz\l{}onem Ginzburga--Landaua w~$\Zp$-parzystym z\l{}o\.zonym
$\hat{\Phi}_i \equiv \hat{s}_i^2$; w~granicy ci\k{a}g\l{}ej odpowiada
$\mathcal{F}^{\mathrm{geo}}_{\mathrm{kin}}[\varphi] = (K_{\mathrm{geo}}/2)
\int \varphi^4 (\nabla\varphi)^2\,d^dx$ (dod.~B,
prop.~\ref{prop:substrate-action}) i~znika to\.zsamo\'sciowo na ka\.zdej
konfiguracji z~$\hat{s}_i = 0$, co kodyfikuje zasad\k{e}
,,brak substratu $\Rightarrow$ brak geometrii''. Jest niezmienniczy
wzgl\k{e}dem \textbf{chiralnej inwersji}:
\begin{equation}\label{eq:Z2-chiral}
    \hat{s}_i \;\mapsto\; -\hat{s}_i
    \qquad \forall\, i \in V.
\end{equation}
Ta dyskretna symetria $\Zp$ nie jest pochodn\k{a} cech\k{a}
hamiltonianu --- jest \textbf{aktem ustanowienia} samego substratu
(patrz Aksjomat~0, dod.~\ref{app:aksjomatyka-roznicy},
tw.~\ref{thm:0-rownowaznosc}). Funkcja stanu $\Ffield$ mierzy
stopie\'n jej spontanicznego z\l{}amania.

\medskip
\noindent\textbf{Etykieta \texttt{prop:Z2}} jest zachowana dla
kompatybilno\'sci wstecznej odno\'snik\'ow; kanoniczna etykieta to
\texttt{ax:Z2-konstytutywny}.
\end{axiom}

\begin{corollary}[Relacja $\Ffield \leftrightarrow
  \langle\hat{s}\rangle$]\label{cor:F-s}
W~przybli\.zeniu pola \'sredniego:
\begin{equation}\label{eq:F-s-relation}
    \Ffield_i[\Omega]
    = \langle\Omega|\hat{s}_i|\Omega\rangle^2.
\end{equation}
$\Ffield = 0 \;\Leftrightarrow\;
\langle\hat{s}_i\rangle = 0$
(symetria $\Zp$ nienaruszona: brak struktury,
nico\'s\'c). $\Ffield > 0 \;\Leftrightarrow\;$
substrat ,,wybra\l{} stron\k{e}'' --- spontaniczne
z\l{}amanie $\Zp$ jednocze\'snie generuje materi\k{e}
i~przestrze\'n.
\end{corollary}

% ----------------------------------------------------------
\subsection{Niestabilno\'s\'c $\Nzero$}\label{ssec:niestabilnosc}

\begin{theorem}[Niestabilno\'s\'c nico\'sci]%
\label{thm:P3}\label{ax:P3}
Stan $\Ffield = 0$ ($\Nzero$) jest niestabilny. Jego miara
w~przestrzeni konfiguracji wynosi zero:
\begin{equation}
    \mu\bigl(\{\sigma \in \Omega : \Ffield[\sigma] = 0\}\bigr) = 0.
\end{equation}
Ka\.zde odej\'scie od dok\l{}adnej symetrii prowadzi do $\Ffield > 0$
lokalnie, co oznacza jednoczesne pojawienie si\k{e} materii
i~generowanej przez ni\k{a} przestrzeni.
\end{theorem}

\begin{proof}[Szkic dowodu (z~Aksjomatu~0)]
Stan $\Ffield = 0$ odpowiada $\langle\hat{s}_i\rangle = 0$ dla
ka\.zdego $i$, czyli konfiguracji \emph{niezmienniczej} wzgl\k{e}dem
$\Zp$. Zgodnie z~prop.~\ref{prop:0-niestabilnosc}
(dod.~\ref{app:aksjomatyka-roznicy}), degeneracja typu
$\sigma x = x$ dla ka\.zdego $x$ jest dynamicznie nieodr\'o\.znialna
od $\Nzero$, st\k{a}d miara konfiguracji $\Zp$-niezmienniczych
jest zero w~$(\Omega, \mu)$. \qed
\end{proof}

\begin{remark}[Status epistemiczny]\label{rem:P3-status}
Dotychczas tw.~\ref{thm:P3} by\l{}o prezentowane jako
\emph{aksjomat}~(A4) substratu. Po promocji $\Zp$ do Aksjomatu~0
staje si\k{e} ono \emph{twierdzeniem} --- konsekwencj\k{a}
logicznej niestabilno\'sci stanu maksymalnie zdegenerowanego.
Etykieta \texttt{ax:P3} jest zachowana dla kompatybilno\'sci
wstecznej.
\end{remark}

Niestabilno\'s\'c nie jest procesem czasowym --- czas nie istnieje
w~$\Nzero$. Jest to stwierdzenie o~\textbf{strukturze} przestrzeni
konfiguracji: stan idealnie symetryczny jest punktem miary zero,
otoczonym ze wszystkich stron stanami z~$\Ffield > 0$.

\begin{proposition}[Niestabilno\'s\'c z~termodynamiki substratu]\label{prop:P3-deriv}
Niech substrat relacyjny $\Gamma = (V, E)$ posiada hamiltoniano
$H_\Gamma$ z~chiralną symetri\k{a} $\Zp$ (inwersja amplitud
$\hat{s}_i \mapsto -\hat{s}_i$). Niech $\mu_\beta$ b\k{e}dzie
miar\k{a} Gibbsa $\mu_\beta(\sigma) \propto e^{-\beta H_\Gamma[\sigma]}$.
W~granicy termodynamicznej ($|V| \to \infty$) istnieje
$\beta_c > 0$ taki, \.ze:
\begin{equation}
    \beta > \beta_c
    \quad\Longrightarrow\quad
    \mu_\beta\bigl(\{\sigma: \Ffield[\sigma] = 0\}\bigr) = 0,
    \qquad
    \mu_\beta\bigl(\{\sigma: \Ffield[\sigma] > 0\}\bigr) = 1.
\end{equation}
Substrat fizyczny (bez zewn\k{e}trznego termostatu) odpowiada
$\beta \to \infty$ --- zawsze jest w~re\.zimie spontanicznie
z\l{}amanej symetrii. Zatem aksjomat~\ref{ax:P3} jest
\textbf{konsekwencj\k{a}} termodynamiki substratu,
nie niezale\.znym postulatem.
\end{proposition}

Analogia: nie pytamy, co ,,wyzwala'' topnienie lodu na ciep\l{}ym blacie.
Pytamy, w~jakich warunkach l\'od jest stabilny. W~TGP: substrat
jest ,,zawsze'' w~warunkach, w~kt\'orych $\Ffield > 0$
jest preferowane --- stan $\Ffield = 0$ ma zerow\k{a} wag\k{e}
statystyczn\k{a}.

\begin{proposition}[Niestabilno\'s\'c $\Nzero$
  z~funkcjona\l{}u dzia\l{}ania]\label{prop:P3-action}
Stan $\Phi = 0$ (nicość) jest niestabilny r\'ownie\.z
w~sensie r\'ownania pola~\eqref{eq:full-field-alt}
(lub~\eqref{eq:covariant-field}). Linearyzuj\k{a}c
r\'ownanie pola wok\'o\l{} $\Phi = 0$:
\begin{equation}\label{eq:P3-lin}
    \nabla^2(\delta\Phi) + \alpha\,\frac{(\nabla\delta\Phi)^2}
    {\delta\Phi} = -q\,\PhiZero\,\rho.
\end{equation}
Ka\.zde niezerowe \'zr\'od\l{}o $\rho > 0$ wymusza
$\delta\Phi > 0$ w~otoczeniu --- pole ,,ucieka'' od zera.
Ponadto potencja\l{} efektywny
$V_{\mathrm{eff}}(\Phi) = (\beta/3)\Phi^3/\PhiZero
- (\gamma/4)\Phi^4/\PhiZero^2$%
\footnote{W~sformu\l{}owaniu kanonicznym TGP (Formulacja~A) potencja\l{} akcji
przyjmuje posta\'c $P(g) = (\beta/7)\,g^7 - (\gamma/8)\,g^8$, gdzie $g = \Phi/\PhiZero$,
a~potencja\l{} statyczny $V(g) = (\gamma/3)\,g^3 - (\gamma/4)\,g^4$.
Cz\l{}on $(\beta/3)\,\psi^3$ jest potencja\l{}em statycznym Formulacji~B.
Patrz: \texttt{tgp\_companion.tex}.}
w~otoczeniu $\Phi = 0$
zachowuje si\k{e} jak $\sim\beta\Phi^3/(3\PhiZero)$, co daje:
\begin{equation}
    V'_{\mathrm{eff}}(0) = 0, \qquad
    V''_{\mathrm{eff}}(0) = 0.
\end{equation}
Brak drugiej pochodnej (masy efektywnej) w~$\Phi = 0$
oznacza, \.ze $\Phi = 0$ jest \textbf{punktem przegi\k{e}cia}
potencja\l{}u, nie minimum --- ka\.zda perturbacja
odsuwa pole od zera. Jest to niezale\.zny argument za
aksjomatem~\ref{ax:P3}, wynikaj\k{a}cy z~samej
struktury r\'ownania pola, nie z~termodynamiki substratu.
\end{proposition}

% ----------------------------------------------------------
\subsection{Dwa sektory}\label{ssec:sektory}

\begin{definition}[Sektory TGP]\label{def:sektory}
Przestrze\'n konfiguracji $\Omega$ dzieli si\k{e} na dwa sektory:
\begin{description}[style=nextline,nosep]
  \item[{$\mathcal{S}_0 = \{\sigma : \Ffield[\sigma] = 0\}$
    --- sektor nico\'sci}]
    Brak materii, brak przestrzeni, brak czasu.
    Stan $\Nzero$ nale\.zy do tego sektora.
  \item[{$\mathcal{S}_1 = \{\sigma : \Ffield[\sigma] > 0\}$
    --- sektor istnienia}]
    Materia obecna, przestrze\'n wygenerowana, czas istnieje
    jako relacja mi\k{e}dzy r\'o\.znymi stanami.
\end{description}
\end{definition}

Relacja $\mathcal{S}_0 \leftrightarrow \mathcal{S}_1$ nie jest
ewolucj\k{a} czasow\k{a} (czas nie istnieje w~$\mathcal{S}_0$).
Jest relacj\k{a} mi\k{e}dzysektorow\k{a} --- dwa r\'o\.zne opisy
tej samej funkcji stanu. Obserwator wewn\k{e}trzny (w~$\mathcal{S}_1$)
widzi bogaty Wszech\'swiat; obserwator zewn\k{e}trzny widzi $\Ffield = 0$
--- nico\'s\'c.

\begin{remark}[Relacja mi\k{e}dzysektorowa jako selekcja miar\k{a}]
,,Przej\'scie'' $\mathcal{S}_0 \to \mathcal{S}_1$ nie zachodzi
\emph{w~czasie} --- czas nie istnieje w~$\mathcal{S}_0$.
Jest to \textbf{selekcja sektora przez miar\k{e}}: prawie
ka\.zda konfiguracja substratu (w~sensie miary $\mu_\beta$
z~stwierdzenia~\ref{prop:P3-deriv}) nale\.zy
do~$\mathcal{S}_1$. Nie ma ,,chwili przed'' ---
sektor $\mathcal{S}_1$ ma pe\l{}n\k{a} miar\k{e}, a~przej\'scie
jest stwierdzeniem o~dominacji statystycznej, nie o~procesie.

To jest fundamentalnie r\'o\.zne od ,,Wielkiego Wybuchu
z~nico\'sci'' --- TGP nie potrzebuje momentu kreacji.
Potrzebuje jedynie faktu, \.ze $\mathcal{S}_0$ jest zerow\k{a}
miar\k{a} w~pe\l{}nej przestrzeni konfiguracji.
\end{remark}

% ----------------------------------------------------------
\subsection{Czas jako konsekwencja $\Ffield > 0$}\label{ssec:czas}

\begin{proposition}[Emergencja czasu]\label{prop:czas}
Czas nie jest postulowany. Wynika z~istnienia wi\k{e}cej ni\.z jednego
stanu:
\begin{equation}
    \Ffield > 0
    \quad\Longrightarrow\quad
    |\{\text{stany}\}| > 1
    \quad\Longrightarrow\quad
    \text{czas istnieje jako relacja mi\k{e}dzy stanami.}
\end{equation}
W~$\Nzero$: $\Ffield = 0$, jeden stan, czas jest poj\k{e}ciowo
zbyteczny.
\end{proposition}

Strza\l{}ka czasu jest wyznaczona przez kierunek wzrostu
entropii relacyjnej $\Srel$ (nie przez wzrost $\Ffield$).

\begin{definition}[Entropia relacyjna $\Srel$]\label{def:Srel}
Niech $\rho_{\mathrm{cg}}(x)$ b\k{e}dzie grubozarnion\k{a}
g\k{e}sto\'sci\k{a} substratu --- miar\k{a} tego, jak
stopnie swobody substratu wype\l{}niaj\k{a} kom\'ork\k{e}
grubozarnienia wok\'o\l{} punktu~$x$, wa\.zon\k{a}
warto\'sci\k{a} $\Phi(x)$. Entropia relacyjna na hypersurface
$\Sigma$ jest:
\begin{equation}\label{eq:Srel}
    \Srel[\Phi] = -\int_\Sigma \rho_{\mathrm{cg}}(x)\,
    \ln\rho_{\mathrm{cg}}(x)\;\sqrt{h}\,d^3x.
\end{equation}
$\Srel$ mierzy r\'o\.znorodno\'s\'c konfiguracji substratu:
im bardziej jednorodna $\rho_{\mathrm{cg}}$, tym wy\.zsza
entropia. $\Srel$ ro\'snie w~trakcie ewolucji (drugie prawo TGP),
co wyznacza strza\l{}k\k{e} czasu.
\end{definition}

\begin{remark}[Niska entropia pocz\k{a}tku]
Spontaniczne z\l{}amanie symetrii (przej\'scie
$\Ffield = 0 \to \Ffield > 0$) generuje stan wysoce
uporz\k{a}dkowany --- jednorodnemu $\Phi$ odpowiada
niska~$\Srel$. P\'o\'zniejsza ewolucja (formowanie struktur,
defekt\'ow, relaksacja gradient\'ow) zwi\k{e}ksza $\Srel$.
W~standardowej kosmologii niska entropia pocz\k{a}tku Wszech\'swiata
jest \emph{zagadk\k{a}} (Past Hypothesis Penrose'a).
W~TGP jest \emph{naturaln\k{a} konsekwencj\k{a}} przej\'scia
fazowego.
\end{remark}

\begin{proposition}[Pr\k{a}d funkcji stanu]\label{prop:F-current}
Funkcja stanu $\Ffield$ posiada lokaln\k{a} posta\'c zachowawcz\k{a}.
Definiujemy \textbf{g\k{e}sto\'s\'c funkcji stanu}
$\mathfrak{f}(x) = \Floc(x) = (\Phi(x) - \PhiZero)^2$ oraz
\textbf{pr\k{a}d funkcji stanu}:
\begin{equation}\label{eq:F-current}
    j_{\Ffield}^\mu(x) \;=\; 2\,(\Phi - \PhiZero)\,\partial^\mu\Phi.
\end{equation}
Wtedy z~r\'ownania pola~\eqref{eq:full-field-alt} wynika:
\begin{equation}\label{eq:F-conservation}
    \partial_\mu j_{\Ffield}^\mu
    \;=\; 2\,(\partial_\mu\Phi)(\partial^\mu\Phi)
    + 2\,(\Phi - \PhiZero)\,\nabla^2\Phi
    \;=\; 2\,(\nabla\Phi)^2
    + 2\,(\Phi - \PhiZero)\bigl[-q\PhiZero\rho
    - \mathcal{N}[\Phi]\bigr].
\end{equation}
Zasada zerowej sumy (aksjomat~\ref{ax:zero}) wymaga, by
ca\l{}ka globalna $\int_\Sigma j_{\Ffield}^0\,\sqrt{h}\,d^3x$
by\l{}a zachowana: nadwy\.zka $\Ffield$ w~jednym regionie jest
kompensowana deficytem w~innym.
W~nieobecno\'sci \'zr\'ode\l{} ($\rho = 0$) i~w~granicy liniowej
($\Phi \approx \PhiZero$, $\mathcal{N} \approx 0$) pr\k{a}d
redukuje si\k{e} do $\partial_\mu j_{\Ffield}^\mu \approx 0$
--- \textbf{\'scis\l{}e prawo zachowania} funkcji stanu
w~pr\'o\.zni.
\end{proposition}

\begin{corollary}[Bilans \'zr\'od\l{}owy funkcji stanu]%
\label{cor:F-source-balance}
Z~r\'ownania~\eqref{eq:F-conservation} wynika, \.ze \'zr\'od\l{}o
materii ($\rho > 0$) stanowi \textbf{ujemny cz\l{}on} po prawej
stronie r\'ownania ci\k{a}g\l{}o\'sci $\Ffield$:
\begin{equation}\label{eq:F-source}
  \partial_\mu j_{\Ffield}^\mu
  = 2\,(\nabla\Phi)^2
  - 2\,q\,\PhiZero\,(\Phi - \PhiZero)\,\rho
  - 2\,(\Phi - \PhiZero)\,\mathcal{N}[\Phi].
\end{equation}
Interpretacja:
\begin{enumerate}[nosep]
  \item \textbf{Cz\l{}on gradientowy} $2(\nabla\Phi)^2 \geq 0$:
    niejednorodne konfiguracje $\Phi$ produkuj\k{a} $\Ffield$ ---
    ka\.zde \emph{istnienie} (struktura) wzmacnia odej\'scie
    od nico\'sci.
  \item \textbf{Cz\l{}on \'zr\'od\l{}owy} $-2q\PhiZero(\Phi-\PhiZero)\rho$:
    w~sektorze $\Phi > \PhiZero$ (wewn\k{a}trz \'zr\'ode\l{})
    materia \emph{absorbuje} $\Ffield$, w~sektorze
    $\Phi < \PhiZero$ (daleko od \'zr\'ode\l{}) materia
    \emph{generuje} $\Ffield$. Ten bilans realizuje
    zasad\k{e} zerowej sumy (ZS2).
  \item \textbf{Cz\l{}on nieliniowy} $-2(\Phi-\PhiZero)\mathcal{N}[\Phi]$:
    samointerferencja pola odpowiada za redystrybucj\k{e} $\Ffield$
    mi\k{e}dzy re\.zimami (przep\l{}yw ,,istnienia'' mi\k{e}dzy skalami).
\end{enumerate}
\end{corollary}

\begin{hypothesis}[Drugie prawo TGP]\label{hyp:second-law}
W~fazie $\Ffield > 0$ r\'ownania efektywne
(r\'ownanie pola~$\Phi$ + kosmologia) implikuj\k{a}:
\begin{equation}\label{eq:second-law}
    \frac{d\Srel}{dt} \geq 0
\end{equation}
dla rozwi\k{a}za\'n kosmologicznych z~warunkiem
pocz\k{a}tkowym odpowiadaj\k{a}cym \'swie\.zemu
przej\'sciu fazowemu ($\Phi \approx \PhiZero$ jednorodnie).

Argument heurystyczny:
\begin{enumerate}[nosep]
  \item Z\l{}amanie symetrii: $\Phi \approx \PhiZero$ jednorodnie
    --- niska $\Srel$ (wysoce uporz\k{a}dkowany stan).
  \item Niestabilno\'s\'c liniowa: perturbacje $\delta\Phi$
    rosn\k{a}, formuj\k{a} si\k{e} struktury --- wzrost $\Srel$.
  \item P\'o\'zna relaksacja: gradienty $\Phi$ si\k{e}
    wyg\l{}adzaj\k{a} --- dalszy wzrost $\Srel$.
\end{enumerate}
\end{hypothesis}

% ----------------------------------------------------------
\subsection{Dynamika substratu poni\.zej progu
  metrycznego}\label{ssec:S0-dyn}

W~$\mathcal{S}_0$ nie istnieje czas w~sensie metrycznym. Jednak
substrat $\Gamma$ posiada w\l{}asn\k{a} mikrodynamik\k{e}, kt\'or\k{a}
mo\.zna opisa\'c parametrem $\tau$ porz\k{a}dkuj\k{a}cym
\'scie\.zki w~przestrzeni konfiguracji --- nie jest to czas,
lecz \textbf{parametr miary na \'scie\.zkach}.

\begin{definition}[Relaksacja stochastyczna substratu]\label{def:relax}
Dynamika substratu w~$\mathcal{S}_0$:
\begin{equation}\label{eq:relax}
    \frac{d\bar{s}_i}{d\tau}
    = -\frac{\partial H_\Gamma}{\partial \bar{s}_i}
    + \eta_i(\tau),
    \qquad
    \langle\eta_i(\tau)\eta_j(\tau')\rangle
    = 2T_{\mathrm{sub}}\,\delta_{ij}\,\delta(\tau-\tau'),
\end{equation}
gdzie $T_{\mathrm{sub}}$ jest \textbf{parametrem
fluktuacji mikrodynamiki} (nie temperatur\k{a}
termodynamiczn\k{a} --- ta wymaga czasu i~geometrii).
\end{definition}

\begin{remark}[Interpretacja $T_{\mathrm{sub}}$]
Fizyczny substrat (bez zewn\k{e}trznego termostatu)
odpowiada $T_{\mathrm{sub}} = 0$, ale kwantowe
relacje komutacyjne
$[\hat{s}_i, \hat{\pi}_j] = i\hbar\delta_{ij}$
generuj\k{a} nieredukowalne fluktuacje. Nawet przy
$T_{\mathrm{sub}} = 0$ stan w~pe\l{}ni symetryczny
jest niestabilny (jak fa\l{}szywa pr\'o\.znia) ---
co jest sp\'ojne z~aksjomatem~\ref{ax:P3}
i~stwierdzeniem~\ref{prop:P3-deriv}.
\end{remark}

\begin{proposition}[Most $\tau \to t$: emergencja czasu metrycznego]%
\label{prop:tau-to-t}
Parametr relaksacyjny $\tau$ substratu (r\'ow.~\ref{eq:relax})
i~czas fizyczny $t$ (wyst\k{e}puj\k{a}cy w~kowariantnym r\'ownaniu
pola) s\k{a} zwi\k{a}zane przez
\textbf{pr\k{e}dko\'s\'c propagacji} $c(\Phi)$:
\begin{equation}\label{eq:tau-to-t}
  dt \;=\; \frac{c(\Phi)}{c_0}\;d\tau
  \;=\; \sqrt{\frac{\PhiZero}{\Phi}}\;d\tau.
\end{equation}
Interpretacja:
\begin{enumerate}[nosep]
  \item W~fazie uporz\k{a}dkowanej ($\Phi = \PhiZero$):
    $dt = d\tau$ --- parametr substratu \textbf{jest} czasem
    fizycznym. Czas metryczny to efektywny opis dynamiki
    substratu wi\k{e}c nie jest dodatkowym stopniem swobody.
  \item W~silnym polu ($\Phi \gg \PhiZero$, np.\ blisko
    \'zr\'od\l{}a): $dt \ll d\tau$ --- czas fizyczny
    \textbf{zwalnia} wzgl\k{e}dem substratu (odpowiednik
    dylatacji grawitacyjnej w~OTW).
  \item W~$\mathcal{S}_0$ ($\Phi = 0$): $dt/d\tau = 0$ ---
    czas fizyczny \textbf{nie p\l{}ynie}, mimo \.ze substrat
    zachowuje mikrodynamik\k{e} (parametr $\tau$ nadal
    istnieje, ale nie generuje obserwowalnych zmian
    metrycznych). To jest \'scis\l{}a realizacja
    aksjomatu~\ref{ax:P3}.
\end{enumerate}
\end{proposition}

\begin{remark}[Strzałka czasu z~termodynamiki substratu]%
\label{rem:arrow-from-substrate}
Relacja~\eqref{eq:tau-to-t} implikuje, \.ze strza\l{}ka czasu
(hipoteza~\ref{hyp:second-law}) jest \textbf{odziedziczona}
po mikrodynamice substratu: relaksacja stochastyczna
(r\'ow.~\ref{eq:relax}) posiada naturalny kierunek
$\tau$ (ku r\'ownowadze termodynamicznej), kt\'ory
przekszta\l{}ca si\k{e} w~fizyczn\k{a} strza\l{}k\k{e}
czasu $t$ przez~\eqref{eq:tau-to-t}. Niska entropia
pocz\k{a}tku Wszech\'swiata wynika z~tego, \.ze
przej\'scie $\mathcal{S}_0 \to \mathcal{S}_1$
tworzy stan daleki od r\'ownowagi substratowej.
\end{remark}

% ----------------------------------------------------------
\subsection{Przej\'scie fazowe $\Nzero \to \mathcal{S}_1$:
  formalizacja Ginzburga--Landaua}\label{ssec:GL-transition}
% ----------------------------------------------------------

Dotychczasowy opis przej\'scia $\mathcal{S}_0 \to \mathcal{S}_1$
(ssec.~\ref{ssec:niestabilnosc}--\ref{ssec:S0-dyn})
opiera si\k{e} na argumentach miarowych i~termodynamicznych.
Poni\.zej formalizujemy je w~ramach \textbf{efektywnej teorii
Ginzburga--Landaua (GL)} dla substratu, pokazuj\k{a}c, jak
r\'ownanie pola~$\Phi$ wynika z~grubowania (coarse-graining)
mikrodynamiki substratu.

\begin{definition}[Funkcjona\l{} GL substratu]\label{def:GL-substrate}
Niech $\phi(x)$ b\k{e}dzie polem porz\k{a}dku substratu
(grubowane $\bar{s}_i$, ci\k{a}g\l{}y odpowiednik zmiennych
substratu, def.~\ref{def:GL-substrate}).
Efektywny funkcjona\l{} Ginzburga--Landaua:
\begin{equation}\label{eq:GL-functional}
  \boxed{
    \mathcal{F}_{\mathrm{GL}}[\phi]
    = \int d^3x\;\biggl[
      \frac{J_{\mathrm{eff}}}{2}\,(\nabla\phi)^2
      + \frac{m_0^2}{2}\,\phi^2
      + \frac{\lambda_0}{4}\,\phi^4
    \biggr],
  }
\end{equation}
z~symetri\k{a} $\Zp$: $\phi(x) \to -\phi(x)$,
gdzie $J_{\mathrm{eff}}$ jest efektywnym sprz\k{e}\.zeniem
gradientowym (proporcjonalnym do $J$ z~$H_\Gamma$),
$m_0^2 = m_0^2(T_{\mathrm{sub}})$ zmienia znak
przy $T_{\mathrm{sub}} = T_c$, a~$\lambda_0 > 0$
stabilizuje potencja\l{}.
\end{definition}

\begin{proposition}[Z\l{}amanie symetrii i~identyfikacja $\Phi$]%
\label{prop:GL-breaking}
Dla $m_0^2 < 0$ (co odpowiada $T_{\mathrm{sub}} < T_c$,
tj.\ fizycznemu substratowi ze~stwierdzenia~\ref{prop:P3-deriv})
minimalizacja $\mathcal{F}_{\mathrm{GL}}$ daje:
\begin{equation}\label{eq:GL-vev}
  \langle\phi\rangle = \pm v_0, \qquad
  v_0 = \sqrt{-m_0^2/\lambda_0}.
\end{equation}
Pole TGP $\Phi$ identyfikujemy z~\textbf{kwadratem
parametru porz\k{a}dku}:
\begin{equation}\label{eq:Phi-from-phi}
  \boxed{
    \Phi(x) = \frac{\phi(x)^2}{\phi_{\mathrm{ref}}^2}\,\PhiZero,
  }
\end{equation}
gdzie $\phi_{\mathrm{ref}}$ jest skal\k{a} odniesienia.
Ta identyfikacja zapewnia:
\begin{enumerate}[nosep]
  \item $\Phi \geq 0$ zawsze (aksjomat~\ref{ax:przestrzen}),
  \item $\Phi = 0$ $\Leftrightarrow$ $\phi = 0$
    (zachowanie $\Zp$, stan $\Nzero$),
  \item $\Phi = \PhiZero$ $\Leftrightarrow$
    $\phi^2 = \phi_{\mathrm{ref}}^2$ (pr\'o\.znia fizyczna),
  \item Symetria $\Zp$ substratu ($\phi \to -\phi$)
    jest \textbf{automatycznie zachowana} na poziomie $\Phi$,
    zgodnie z~aksjomatem~\ref{ax:zero} (zasada zerowej sumy ZS1).
\end{enumerate}
\end{proposition}

\begin{proof}
Wystarczy sprawdzi\'c warunki krytyczne.
(1)~$\Phi = \phi^2/\phi_{\mathrm{ref}}^2 \cdot \PhiZero \geq 0$
--- trywialnie, bo $\phi^2 \geq 0$.
(2)~Z~$\Phi = 0$ wynika $\phi = 0$ (i~odwrotnie).
(3)~Z~$\Phi = \PhiZero$ wynika $\phi^2 = \phi_{\mathrm{ref}}^2$,
wi\k{e}c $\phi = \pm\phi_{\mathrm{ref}}$.
(4)~$\phi \to -\phi$ daje $\Phi \to \Phi$ ---
$\Phi$ jest $\Zp$-niezmiennicze.
\end{proof}

\begin{remark}[SSB wymaga aksjomatycznie $m_0^2 < 0$ (v2) --- 2026-04-25]
\label{rem:GL-breaking-axiomatic}
W~v2-aksjomatyce wi\k{a}zanie gradientowe Ginzburga--Landaua
(\eqref{eq:H-Gamma} w~dod.~\ref{app:substrat},
rem.~\ref{rem:B-ssb-v2}) znika na dowolnej konfiguracji
jednorodnej, wi\k{e}c \emph{nie wnosi wk\l{}adu do
\'sredniopolowego $v^2$}. Inaczej ni\.z w~wersji v1
(bilinear\-ne~wi\k{a}zanie $-J\sum \hat{s}_i\hat{s}_j$
z~alternatywnym kryterium $Jz > m_0^2$), jedynym mechanizmem
SSB w~v2 jest $m_0^2(T_{\mathrm{sub}}) < 0$. Dlatego warunek
stwierdzenia~\ref{prop:GL-breaking} nie jest twierdzeniem
wyprowadzonym z~dynamiki wi\k{a}za\'n, lecz \emph{konsekwencj\k{a}}
aksjomatu~\ref{ax:substrat} (A5: $m_0^2$ zmienia znak
przy~$T_c$). Istnienie $T_c$ jest zatem \emph{danymi
aksjoma\-tycz\-nymi}, nie rezul\-ta\-tem dynamiki.
\end{remark}

\begin{proposition}[Wyprowadzenie r\'ownania pola $\Phi$
  z~GL]\label{prop:field-from-GL}
Przepisuj\k{a}c $\mathcal{F}_{\mathrm{GL}}$ w~zmiennej
$\Phi = \phi^2\,\PhiZero/\phi_{\mathrm{ref}}^2$
(zak\l{}adaj\k{a}c $\phi > 0$ w~jednej domenie) i~variuj\k{a}c
$\delta\mathcal{F}/\delta\Phi = 0$, otrzymujemy:
\begin{equation}\label{eq:GL-to-TGP}
  \frac{J_{\mathrm{eff}}}{4\PhiZero}\!\left[
    \nabla^2\Phi + 2\,\frac{(\nabla\Phi)^2}{\Phi}
  \right]
  + \frac{m_0^2}{2}\,\sqrt{\frac{\Phi}{\PhiZero}}
  + \lambda_0\,\frac{\Phi}{\PhiZero}
  = 0.
\end{equation}
W~przybli\.zeniu $\Phi \approx \PhiZero$ (s\l{}abe pole),
rozwijaj\k{a}c $\sqrt{\Phi/\PhiZero}
= \sqrt\psi \approx 1 + \delta\Phi/(2\PhiZero)$
i~identyfikuj\k{a}c:
\begin{equation}\label{eq:GL-identification}
  \alpha = 2, \qquad
  \beta = \frac{2\lambda_0\,\PhiZero}{J_{\mathrm{eff}}}, \qquad
  \gamma = \frac{2\lambda_0\,\PhiZero}{J_{\mathrm{eff}}},
\end{equation}
odzyskujemy r\'ownanie pola TGP~\eqref{eq:full-field-alt}
z~$\beta = \gamma$ --- \textbf{warunek pr\'o\.zniowy wynika
automatycznie} z~symetrii GL!
\end{proposition}

\begin{proof}
Podstawiamy $\phi = \sqrt{\PhiZero^{-1}\,\phi_{\mathrm{ref}}^2\,\Phi}$:
\[
  \nabla\phi = \frac{\phi_{\mathrm{ref}}}{2\sqrt{\PhiZero\Phi}}\,\nabla\Phi,
  \qquad
  (\nabla\phi)^2 = \frac{\phi_{\mathrm{ref}}^2}{4\PhiZero}\,
    \frac{(\nabla\Phi)^2}{\Phi}.
\]
Cz\l{}on gradientowy:
$\frac{J_{\mathrm{eff}}}{2}(\nabla\phi)^2
= \frac{J_{\mathrm{eff}}\phi_{\mathrm{ref}}^2}{8\PhiZero}
  \frac{(\nabla\Phi)^2}{\Phi}$.

Cz\l{}on masowy:
$\frac{m_0^2}{2}\phi^2
= \frac{m_0^2\,\phi_{\mathrm{ref}}^2}{2\PhiZero}\,\Phi$.

Cz\l{}on $\lambda$:
$\frac{\lambda_0}{4}\phi^4
= \frac{\lambda_0\,\phi_{\mathrm{ref}}^4}{4\PhiZero^2}\,\Phi^2$.

Wariacja $\delta\mathcal{F}/\delta\Phi = 0$:
\[
  -\frac{J_{\mathrm{eff}}\phi_{\mathrm{ref}}^2}{8\PhiZero}\!\left[
    \frac{\nabla^2\Phi}{\Phi}
    + 2\,\frac{(\nabla\Phi)^2}{\Phi^2} - \ldots
  \right]
  + \frac{m_0^2\phi_{\mathrm{ref}}^2}{2\PhiZero}
  + \frac{\lambda_0\phi_{\mathrm{ref}}^4}{2\PhiZero^2}\,\Phi
  = 0.
\]
Mno\.z\k{a}c przez $-8\PhiZero\Phi/(J_{\mathrm{eff}}\phi_{\mathrm{ref}}^2)$
i~linearyzuj\k{a}c wok\'o\l{} $\Phi = \PhiZero$, po uporz\k{a}dkowaniu
dostajemy:
\[
  \nabla^2\Phi + 2\,\frac{(\nabla\Phi)^2}{\Phi}
  + \beta\,\frac{\Phi^2}{\PhiZero}
  - \gamma\,\frac{\Phi^3}{\PhiZero^2} = 0
\]
z~$\beta = \gamma = 2\lambda_0\PhiZero/J_{\mathrm{eff}}$,
co jest dok\l{}adnie r\'ownaniem pola~\eqref{eq:full-field-alt}
z~$\alpha = 2$ i~$\beta = \gamma$ (pr\'o\.znia).
Czynnik $\alpha = 2$ wynika z~$\Phi = \phi^2$ ---
sprz\k{e}\.zenie gradientowe jest geometryczne.
\end{proof}

\begin{remark}[Fizyczna tre\'s\'c wyniku]\label{rem:GL-physics}
Stwierdzenie~\ref{prop:field-from-GL} jest kluczowe,
poniewa\.z:
\begin{enumerate}[nosep]
  \item \textbf{Warunek $\beta = \gamma$} nie jest
    niezale\.znym postulatem, lecz \emph{konsekwencj\k{a}}
    symetrii $\Zp$ substratu i~identyfikacji $\Phi = \phi^2$.
  \item \textbf{Eksponent $\alpha = 2$} nie jest wybrany
    ,,z~dzia\l{}ania'' (jak w~hipotezie~\ref{hyp:action}),
    lecz \emph{wynika} z~relacji $\Phi \propto \phi^2$:
    $\nabla\Phi/\Phi = 2\nabla\phi/\phi$,
    sk\k{a}d sprz\k{e}\.zenie gradientowe daje
    \emph{dok\l{}adnie} $\alpha = 2$.
  \item \textbf{Parametry fizyczne} $\beta$, $\gamma$
    s\k{a} wyra\.zone przez sta\l{}e substratu:
    $\beta = \gamma = 2\lambda_0\PhiZero/J_{\mathrm{eff}}$.
    Pomiar $\beta$ (lub $\gamma$) z~danych kosmologicznych
    pozwala ograniczy\'c stosunek $\lambda_0/J_{\mathrm{eff}}$.
  \item \textbf{Przej\'scie fazowe} jest typu \emph{drugiego rz\k{e}du}
    (klasy uniwersalno\'sci 3D Isinga ze skalarem O(1)),
    co jest sp\'ojne z~wynikami renormalizacji substratu
    (punkt Wilsona--Fishera, patrz ANALIZA\_SPOJNOSCI\_v4,
    \S6b).
\end{enumerate}
\end{remark}

\begin{remark}[Defekty topologiczne i~materia]\label{rem:GL-defects}
W~fazie z\l{}amanej symetrii ($\phi \neq 0$) mog\k{a} istnie\'c
\textbf{\'sciany domenowe} --- regiony, w~kt\'orych $\phi$
przechodzi przez zero (zmienia domen\k{e} $+v_0 \to -v_0$).
Na \'scianie $\phi = 0$, wi\k{e}c $\Phi = 0$ ---
lokalna ,,nico\'s\'c'' osadzona w~istniej\k{a}cej przestrzeni.
W~trzech wymiarach:
\begin{itemize}[nosep]
  \item \'Sciany domenowe (dim = 2): niestabilne, kurczą się
    (energia $\propto$ pole powierzchni).
  \item Struny (dim = 1): liniowe defekty; w~TGP potencjalnie
    odpowiadaj\k{a} strunom kosmicznym.
  \item Konfiguracje je\.za (dim = 0): lokalne $\phi = 0$
    z~topologicznym nawinięciem; kandydaci na cz\k{a}stki
    (hipoteza~\ref{hyp:spin-chirality} w~sek.~\ref{ssec:fermiony}).
\end{itemize}
Relacja mi\k{e}dzy defektami topologicznymi GL
a~fundamentaln\k{a} materi\k{a} TGP pozostaje
\textbf{otwartym problemem}.
\end{remark}

% ----------------------------------------------------------
\subsection{S\l{}ownik interpretacyjny}\label{ssec:slownik-interp}

\begin{remark}[Interpretacja podstawowych byt\'ow TGP]\label{rem:interpretacja}
\begin{description}[style=nextline, nosep]
  \item[\textbf{Materia}] Stabilna klasa konfiguracji substratu
    ze z\l{}aman\k{a} symetri\k{a} $\Zp$. \'Zr\'od\l{}o pola $\Phi$.
    Nie ,,mieszka w~przestrzeni'' --- \emph{generuje} przestrze\'n.
  \item[\textbf{Przestrze\'n}] Faza metryczna substratu, opisana
    polem $\Phi > 0$. Nie jest aren\k{a} --- jest \emph{produktem}
    materii. Istnieje tylko tam, gdzie $\Phi > 0$.
  \item[\textbf{Pr\'o\.znia (lokalna)}] Region, w~kt\'orym brak
    lokalnych \'zr\'ode\l{}, ale $\Phi \approx \PhiZero > 0$.
    Przestrze\'n istnieje --- jest generowana przez globaln\k{a}
    materi\k{e}.
  \item[\textbf{Nico\'s\'c ($\Nzero$)}] Brak przestrzeni ($\Phi = 0$),
    materii, czasu. Fundamentalnie r\'o\.zna od pr\'o\.zni.
  \item[\textbf{Pole $\Phi$}] Parametr fazy metrycznej substratu.
    Mierzy stopie\'n, w~jakim substrat wspiera relacje metryczne
    (odleg\l{}o\'sci, propagacj\k{e}). Nie jest ,,polem na
    przestrzeni'' --- jest polem, \emph{z~kt\'orego} przestrze\'n
    si\k{e} sk\l{}ada.
  \item[\textbf{Geometria ($g_{\mu\nu}$)}] Obiekt emergentny:
    efektywny opis konfiguracji $\Phi$ na skalach, gdzie pole
    jest g\l{}adkie. Nigdy nie jest fundamentem --- zawsze jest
    \emph{konsekwencj\k{a}} pola $\Phi$.
  \item[\textbf{Czas}] Relacja porz\k{a}dkuj\k{a}ca stany
    w~$\mathcal{S}_1$. Nie istnieje w~$\Nzero$.
    Strza\l{}ka czasu wynika ze wzrostu $\Srel$.
  \item[\textbf{\'Zr\'od\l{}o}] Dowolna konfiguracja substratu
    generuj\k{a}ca wk\l{}ad do $\Phi$: cz\k{a}stki, pola, energia.
  \item[\textbf{Defekt}] Zlokalizowany, topologicznie stabilny
    profil $\Phi(r) \neq \PhiZero$ --- kandydat na cz\k{a}stk\k{e}.
  \item[\textbf{Energia}] W\l{}asno\'s\'c dynamiczna \'zr\'od\l{}a
    lub konfiguracji pola $\Phi$. Nie jest oddzielnym bytem.
\end{description}
\end{remark}

\begin{remark}[Co jest obserwowalne?]\label{rem:obserwowalne}
\begin{description}[style=nextline, nosep]
  \item[\textbf{Obserwowalne bezpo\'srednio}:]
    $g_{\mu\nu}^{\mathrm{eff}}$ (metryk\k{a} efektywna, dost\k{e}pna
    pomiarom); propagacja foton\'ow i~cz\k{a}stek;
    efektywne sta\l{}e $c_{\mathrm{lok}}, \hbar_{\mathrm{lok}},
    G_{\mathrm{eff}}$; dynamika struktur kosmicznych.
  \item[\textbf{Obserwowalne po\'srednio}:]
    Pole $\Phi$ --- przez swoje efekty na metryk\k{e},
    sta\l{}e i~dynamik\k{e}. Warto\'s\'c $\Phi$ nie jest
    mierzalna bezpo\'srednio, ale stosunki $\Phi_1/\Phi_2$
    mi\k{e}dzy regionami --- tak (przez stosunek $c_1/c_2$
    lub $G_1/G_2$).
  \item[\textbf{Nieobserwowalne bezpo\'srednio}:]
    Substrat $\Gamma = (V, E)$; operatory $\hat{s}_i$;
    symetria $\Zp$ na poziomie mikroskopowym;
    parametr $\tau$ mikrodynamiki substratowej.
    S\k{a} to elementy \emph{modelu ukrytego}, dost\k{e}pne
    jedynie po\'srednio, przez konsekwencje na poziomie $\Phi$.
\end{description}
\end{remark}

% ----------------------------------------------------------
\subsection{Wsp\'o\l{}zale\.zno\'s\'c materii i~przestrzeni}\label{ssec:wspolzaleznosc}

Centralna r\'o\.znica wzgl\k{e}dem standardowej fizyki:

\begin{center}
\begin{tabular}{lcc}
\toprule
& \textbf{Standardowa fizyka} & \textbf{TGP} \\
\midrule
Przestrze\'n & istnieje niezale\.znie & generowana przez materi\k{e} \\
Materia & ,,w'' przestrzeni & \'zr\'od\l{}o przestrzeni \\
Pr\'o\.znia & pusta przestrze\'n & t\l{}o referencyjne $\Phi = \PhiZero > 0$
  (generowane globalnie); $\Phi = 0$ to nico\'s\'c, nie pr\'o\.znia \\
Grawitacja & krzywizna gotowej geometrii & efekt interferencji $\Phi$ \\
Czas & koordynata & relacja w~$\mathcal{S}_1$ \\
\bottomrule
\end{tabular}
\end{center}

Kluczowe: materia i~przestrze\'n nie s\k{a} osobnymi bytami ---
s\k{a} dwoma aspektami tego samego zjawiska. Cz\k{a}stka
\textbf{jest} lokalnym wzbudzeniem, kt\'ore jednocze\'snie
generuje przestrze\'n wok\'o\l{} siebie. Przestrze\'n \textbf{jest}
polem $\Phi$ wytworzonym przez materi\k{e}.
